% TOP_PROBLEM: {"id": 3234, "title": "Counting 3-colorings of the hex lattice", "category_id": 3, "category": {"id": 3, "name": "graph_theory", "display_name": "Graph Theory", "description": "Problems involving graphs, networks, and their properties.", "slug": "graph-theory", "order_index": 3, "created_at": "2026-07-31T15:26:25.671Z"}, "status": "open", "problem_number": "OPG-1764", "set_id": 12}
% FIRST_POSTED: 2026-10-01
% ENTRY_KIND: statement_only
% UnsolvedMath ID 3234; source catalogue code OPG-1764
% !TeX program = lualatex
% Definition and sources only; this file is not an LLM solution attempt.
\documentclass[11pt,a4paper]{article}
\usepackage[margin=25mm]{geometry}
\usepackage{fontspec}
\setmainfont{lmroman10-regular.otf}[BoldFont=lmroman10-bold.otf,
 ItalicFont=lmroman10-italic.otf,BoldItalicFont=lmroman10-bolditalic.otf]
\usepackage{amsmath,amssymb,mathtools}
\usepackage{unicode-math}
\setmathfont{latinmodern-math.otf}
\AtBeginDocument{\let\setminus\smallsetminus}
\usepackage{xurl,hyperref}
\hypersetup{colorlinks=true,urlcolor=blue}
\setcounter{secnumdepth}{0}
\setlength{\parindent}{0pt}
\setlength{\parskip}{5pt}
\setlength{\emergencystretch}{3em}
\begin{document}
\section*{Counting 3-colorings of the hex lattice}\label{3234}
{\small UnsolvedMath ID 3234 · OPG-1764 · Graph Theory · OpenGarden}
\subsection{Definitions and mathematical statement}
Tile the Euclidean plane by unit equilateral triangles. For each integer $n\ge1$, take the regular hexagonal patch whose six sides consist of $n$ triangle edges. Identify each side with its opposite side by translation, matching the subdivision vertices. This quotient is a triangulated torus $T_n$. Keep the resulting cellular map structure, including distinct edges if identifications produce parallel edges at a small value of $n$.

Let $H_n$ be the geometric dual of this torus triangulation: it has one vertex for each triangular face of $T_n$, and an edge across each primal edge joining its incident face vertices. Thus $H_n$ is the toroidal hexagonal lattice map associated with this specific periodic patch. There are $6n^2$ triangular faces before identifications, so $|V(H_n)|=6n^2$.

For any finite graph $G$, write $P_G(3)$ for the number of maps $c:V(G)\to\{1,2,3\}$ assigning different labels to the endpoints of every edge. Color labels are distinguished; colorings related by a global color permutation are counted separately. The problem is to prove convergence and determine the exact value of
\[
                \lim_{n\to\infty}P_{H_n}(3)^{1/(6n^2)}.
\]
This is the exponential growth factor per vertex for vertex colorings of the stated periodic hexagonal graphs. It is not an edge-coloring or face-coloring model. The specified torus identifications are part of the sequence, rather than free boundary conditions on a planar patch.
\subsection{Short English statement}
Find the limiting number of proper three-colorings per vertex for the specified periodic hexagonal lattice on a torus.
\subsection{Sources}
\begin{itemize}
\item[S1] ULAM.AI and the UnsolvedMath Contributors, supplied problems.json, record 3234 (OPG-1764), OpenGarden. Catalogue identity and source transcription; CC BY 4.0.
\url{https://huggingface.co/datasets/ulamai/UnsolvedMath}.
\item[S2] Open Problem Garden, Counting 3-colorings of the hex lattice. Original problem and discussion.
\url{http://www.openproblemgarden.org/op/counting_3_colorings_of_the_hex_lattice}.
\end{itemize}
\end{document}
